% --- CONFIGURAZIONE OTTIMIZZATA PER APPUNTI DI INGEGNERIA DEL SOFTWARE ---

% 1. Lingua e Codifica
\usepackage[T1]{fontenc}
\usepackage[utf8]{inputenc}
\usepackage[italian]{babel}

% 2. Layout e Geometria (Sostituisce tutti i calcoli manuali)
\usepackage[a4paper]{geometry}
\geometry{top=2.5cm, bottom=2.5cm, left=2.5cm, right=2.5cm, headheight=15pt}
\usepackage{parskip} % Gestisce lo spazio tra i paragrafi senza indentazione

% 3. Grafica e Colori
\usepackage{graphicx}
\usepackage{xcolor}
\definecolor{codegreen}{rgb}{0,0.6,0}
\definecolor{codegray}{rgb}{0.5,0.5,0.5}
\definecolor{codepurple}{rgb}{0.58,0,0.82}
\definecolor{backcolour}{rgb}{0.95,0.95,0.92}
\definecolor{unict_blue}{RGB}{0,39,75}

% 4. Codice (Essenziale per Informatica)
\usepackage{listings}
\lstset{
    backgroundcolor=\color{backcolour},   
    commentstyle=\color{codegreen},
    keywordstyle=\color{magenta},
    numberstyle=\tiny\color{codegray},
    stringstyle=\color{codepurple},
    basicstyle=\ttfamily\small,
    breakatwhitespace=false,         
    breaklines=true,                 
    captionpos=b,                    
    keepspaces=true,                 
    numbers=left,                    
    numbersep=5pt,                  
    showspaces=false,                
    showstringspaces=false,
    showtabs=false,                  
    tabsize=2,
    language=SQL % Default per il tuo corso
}

% 5. Matematica e Simboli
\usepackage{amsmath, amssymb, amsthm, bm}
\usepackage{pgfplots}
\pgfplotsset{compat=1.18} % Risolve il warning che vedevi

% 6. Intestazioni e Pié di pagina (Fancyhdr pulito)
\usepackage{fancyhdr}
\pagestyle{fancy}
\fancyhf{} % Pulisce tutto
\fancyhead[L]{\nouppercase{\leftmark}} % Capitolo a sinistra
\fancyhead[R]{\thepage} % Numero pagina a destra
\renewcommand{\headrulewidth}{0.4pt}


% 7. Ambienti per Definizioni (Belli da vedere)
\usepackage[most]{tcolorbox}
\newtcbtheorem[number within=chapter]{definition}{Definizione}{
    enhanced, % <--- FONDAMENTALE PER RENDERIZZARE IL TITOLO ESTERNO
    colback=unict_blue!5,
    colframe=unict_blue,
    coltitle=white, % <--- TESTO BIANCO PER LEGGIBILITA'
    fonttitle=\bfseries,
    separator sign={:},
    boxrule=0.5pt,
    attach boxed title to top left={yshift=-2mm, xshift=2mm},
    boxed title style={colback=unict_blue}
}{def}

% 8. Altri pacchetti utili
\usepackage{enumitem}
\usepackage{hyperref} % Link cliccabili nell'indice
\hypersetup{
    colorlinks=true,
    linkcolor=unict_blue,
    filecolor=magenta,      
    urlcolor=cyan,
    pdftitle={Appunti Ingegneria del Software},
}
\usepackage{tikz}
\usetikzlibrary{positioning}
\usepackage{float}
